\documentclass[11pt]{article}
\usepackage[margin=1in]{geometry}
\usepackage[T1]{fontenc}
\usepackage{lmodern}
\usepackage{microtype}
\usepackage{hyperref}
\usepackage{enumitem}
\hypersetup{colorlinks=true,linkcolor=blue,urlcolor=blue}
\setlength{\emergencystretch}{1.5em}
\setlist[itemize]{itemsep=0.25em,topsep=0.4em}
\title{\textit{cs-mail}: Aliases and Two Forms of Express-Lane Permission}
\author{C-SQD}
\date{Product design memo --- 21 October 2026}
\begin{document}
\maketitle

\paragraph{Status and authority.}
This memo explains the current alias design. The
\href{DOMAIN_MODEL.md}{C-SQD domain model} owns its confirmed product rules.
Alias management, receiving inboxes and lane management now use independent
receiving policy and per-alias publications. The 21 October implementation replaces
the earlier public/private contracts; current validation and owner-install status
are recorded in the Apple runbook. The
\href{cs_mail_build_plan.tex}{build plan} owns sequencing. The
\href{docs/rust-domain-principles.md}{Rust-domain principles} remain binding.

\section{The Problem: One Address, Many Different Relationships}

A person gives an email address to a hospital, an insurer, a pharmacy and many
other organizations. Correspondence from all of them reaches the same entrance.
When the address is circulated, the recipient has little context about which
relationship supplied that route. Replacing it can disrupt unrelated contacts,
portal access and account notices.

Expected healthcare correspondence also comes from people the recipient has not
contacted before. A hospital may pass contact details to a specialist, and its
departments or service vendors may use different sending addresses. Requiring
each correspondent to make a paid approach would put the request process in the
way of an existing organizational relationship.

A publicly listed address serves a different purpose. A professor may welcome
new research proposals while needing control over unsolicited approaches. The
existing bonded request gives a stranger an explicit way to ask for permission.
Both uses belong in CS-MAIL, within the same person's account.

\section{One Alias Model, Two Receiving Policies}

Agreed 21 October 2026. Each member has one protocol identity and multiple aliases.
Public/private is no longer an alias kind, and there is no single-current-public
limit. For any alias the member may choose a readable address or request generation,
independently of receiving policy and directory listing. An owner-visible
purpose/description stays undisclosed unless explicitly shared.

\section{An Alias for a Hospital}

The member selects Allow contact through this address and explicitly authorizes
an alias-scoped express lane. A hospital or a referred specialist can correspond
without a bonded request, subject to authentication, coverage, blocks and admission
controls. This receiving policy has no bonded fallback. Possession establishes
neither organizational affiliation nor who shared the address; a leaked address
can reach other senders. This policy is a member choice, never inferred from the
purpose label or the assumption that healthcare or government is lower risk.

\section{A Bonded Alias for Shopping or Newsletters}

The member selects Require a bonded request. Each such alias owns its profile and
request classes. The member can grant a sender-scoped lane for explicitly selected
authenticated domains/mailboxes or standing permission under a named receiving
scope. Those senders can correspond directly; other eligible senders must select
one of this alias's classes and explicitly authorize a bonded request. A sender
mismatch or ordinary delivery attempt never authorizes a charge. An organization
name never determines which domains are authorized.

\section{A Bonded Alias for Professor Bob}

Bob can choose a readable or generated address and independently select directory
visibility. Each of his bonded aliases has its own profile and classes. An unlisted
bonded alias remains approachable by exact address or direct link; multiple aliases
can serve different purposes without publishing their common ownership. Alice
selects the particular alias's class and reviews the existing conditional charge.
Bob's acceptance, rejection and expiry retain the existing settlement rules.
Blocking cannot be bypassed by paying or choosing another receiving alias.

\section{Creation, Transitions and Future Pricing}

One New alias journey collects purpose, address choice, receiving policy and,
for a bonded alias, its profile/classes and separate visibility choice. Creation
alone grants no lane; the member's configuration explicitly authorizes permission.
Rust owns distinct meanings, current authority, complete effects, atomic commit
and exact retry. Swift preserves edits and place, presenting the member's choices
and verified outcome without implementation machinery. Rust and Swift now implement
these contracts, with an explicit offline preservation upgrade and one current
runtime representation.

Receiving-policy switching needs a defined transition for pending requests,
frozen terms, profiles/classes and existing grants. Address/lane restriction,
removal and retirement never silently switch policy or settle requests. Replacement,
reply addressing and held-message handling retain their explicit open decisions.

A nominal charge per new alias associated with the annual utility-fee cycle is a
future direction. Price, recurrence versus creation-only charging, allowances,
mid-cycle treatment, collection and failed/retried/replaced alias treatment remain
open. No charge is active or authorized by this note. Fees must remain distinct
from bonded request charges, annual service and member distributions.


\section{Two Bases for Express-Lane Permission}

\paragraph{Alias-scoped.}
The member explicitly grants a lane for a named receiving alias. Permission is bound
to that receiving address and its active lane policy. It can accommodate a new
sender who has been given the address without requiring individual approval first.

\paragraph{Sender-scoped.}
The member explicitly authorizes a sender subject under a grant's scope and
validity. Existing lanes may bind to a native identity, an external mailbox or a
recipient-selected legacy domain. The existing
\href{cs_mail_express_lanes.tex}{express-lane memo} explains this model, and its
\href{cs_mail_express_lanes_implementation.tex}{implementation plan} records the
current sender-scoped machinery.

An alias is a communication address and managed route; a lane is permission.
Creation alone grants no lane; the member explicitly selects receiving policy
and permission. A bonded alias can grant direct access to selected senders;
an alias-scoped policy grants its broad lane explicitly. An active address supplies no authority
without an applicable permission. Lane changes do not retire the address, and
route restriction or retirement constrains every permission through it.
Standing acceptance remains a distinct permission. A protocol identity identifies a participant; an alias is a separately identified
route controlled by that participant. One protocol identity may control multiple
aliases. Alias creation, replacement and retirement preserve that identity,
account, login and shared subject. The association between aliases must not
be automatically disclosed to correspondents. This supersedes the old
address-equals-protocol-identity product assumption; implementation awaits cutover.

\section{What the Route Tells the Member}

The hospital alias identifies the entrance the member supplied for that
relationship. It does not prove that every message was authored by the hospital,
that every new sender is its authorized delegate, or that the hospital caused an
address leak. An alias can be forwarded or disclosed. Sender authentication
supplies separate evidence about the message's origin.

The design distinguishes route, evidence and permission. Previously unknown
senders need no individual sender grant only when an applicable alias-scoped
lane authorizes that receiving route. An address without applicable permission is insufficient.
Suspicious or insufficiently authenticated correspondence is held separately for
explicit member review. Holding or releasing one message grants no ongoing
permission and creates no payment invitation. A hold cannot bypass an established
identity-level block; unresolved authentication cannot establish that a sender is
unblocked. Detailed release, retention and refusal contracts remain open.

SMTP establishes the receiving alias from the accepted envelope recipient, not
the visible To header. Native commands authenticate the target alias. Every
admitted message retains its receiving alias and permission basis through storage
and replies; replacement must not rewrite the original route.

External email retains its existing privacy boundaries: receiving and relay
services may handle plaintext. The alias does not confer endpoint-only
confidentiality. Existing disclosure-consent requirements remain applicable;
the alias design does not itself authorize outbound disclosure or managed decryption.

\section{CS-MAIL Owns the Infrastructure; Phoros Uses It Later}

CS-MAIL, operated by C-SQD, PBC, owns alias creation, routing, permissions,
lifecycle management and correspondence. The member must be able to use these
capabilities independently of Phoros.

Phoros Account, operated by Reservatory, PBC, will later use the infrastructure
through a linked CS-MAIL account. It can associate healthcare relationships with
their communication routes without building a second alias system. Product
linking does not itself grant access to messages, authorize decryption or disclose
healthcare records. The \href{docs/product-identity-model.md}{product identity
model} continues to distinguish product accounts, logins and shared subject
association. No primary recovery email address is introduced by this design.

\section{Blocking and Restrict, Replace, Retire}

A sender protocol-identity block overrides every permission path across all
receiving aliases controlled by the recipient. An external mailbox identifies
only what authentication establishes; two mailboxes are not automatically one
person. Aliases have no automatic expiry. Their controls are:
\begin{description}
  \item[Restrict.] Stop new correspondents while preserving authorized existing
  conversations. Only exact conversations with admitted correspondence at restriction qualify; new conversations require explicit approval.
  \item[Replace.] Create a new alias with an explicit transition for the old one.
  Preserve identity, blocks, correspondence and original route history. Do not
  automatically carry unknown-sender admission into the new route.
  \item[Retire.] Stop all inbound delivery through the alias and preserve received
  correspondence. Explain that notices and recovery messages will stop arriving.
  Never reassign the retired address to another member.
\end{description}
None switches receiving policy, creates a paid request or settles an existing one. Replacement
does not update an external portal automatically. A transition must explicitly
address portal login, recovery and notices. Outbound composition identifies the
selected alias; detailed reply continuity still requires design.

\paragraph{Lane actions.}
Restricting a broad lane preserves only exact already admitted conversations.
Removing it withdraws its authority for subsequent messages, including existing
conversations. Replacing it with sender-scoped permission commits withdrawal and
the selected grants atomically, without grandfathering earlier senders. Other
independently valid permission remains separate. These actions preserve the
address, inbox and received copies; they differ from replacing or retiring the
address. Reopening a lane is explicit and never releases refused or held mail.

\section{Decisions Still Required}

The participant/route distinction, identity-level blocking, no automatic alias
expiry, separate lane permission, lifecycle actions and held-message direction
are agreed. Still open are:
\begin{itemize}
  \item Replacement transition duration and portal, recovery and notice continuity.
  \item Detailed outbound selection, reply routing and held-message review/retention.
  \item Detailed scoped identity-reference and held-message contracts.
  \item Receiving-policy switching, including pending requests and existing grants.
  \item Initial directory visibility and alias-fee pricing/collection semantics.
  \item Milestone placement, member-management flows and validation scope.
\end{itemize}
Do not infer all-alias permission from a sender grant. Alias creation must not
inherit the pending-request settlement effects of sender-lane management. Existing
sender-scoped implementation and historical financial outcomes require an explicit
preservation-checked cutover; validation scope is recorded in the Apple runbook.

\end{document}
